\section{Critical bridge length and a polynomial polygon observable}
\label{sec:pv:introduction}

A walk can cross a wide strip while making many returns to a much smaller
region. Its total crossing weight therefore does not determine its mean
length. We study this distinction for critical self-avoiding walks on the
honeycomb lattice. Our initial boundary point is fixed and the terminal
point is summed over the opposite wall. The main result is a mean-length
exponent of $4/3$ under this all-length measure.

We use the honeycomb lattice dual to the equilateral triangular tiling of
side one. Its edges have length $1/\sqrt3$. A \emph{port} is the midpoint
of an edge cut by a tiling line. A path between ports visits distinct
honeycomb vertices, and its length $|\gamma|$ is their number. Its critical
weight is $x_*^{|\gamma|}$, where
\[
 x_*=(2\cos(\pi/8))^{-1}=(2+\sqrt2)^{-1/2}.
\]
The terminal half-edges have no additional weight. A strip of $h$ layers
has physical height $h\sqrt3/2$. With a bottom source at horizontal
coordinate zero, its attainable top ports lie in $\mathbb Z+h/2$.
Write $\mathcal B_h(e)$ for the paths confined to that strip from the
source to top port $e$, and put
\[
 b_h(e)=\sum_{\gamma\in\mathcal B_h(e)}x_*^{|\gamma|},\qquad
 b_h=\sum_e b_h(e),\qquad
 \ell_h=\sum_e\sum_{\gamma\in\mathcal B_h(e)}
 |\gamma|x_*^{|\gamma|}.
\]
Thus $\ell_h$ is a first-length insertion, rather than a partition
function conditioned on one length. The bridge law assigns each such
path probability $x_*^{|\gamma|}/b_h$.

\begin{theorem}[Strip mass and bridge length]\label{thm:pv:bridge}
As the integer layer height $h$ tends to infinity,
\[
 b_h=h^{-1/4+o(1)},\qquad
 \ell_h=h^{13/12+o(1)},\qquad
 \mathbb E_h|\gamma|=h^{4/3+o(1)}.
\]
All three sums are finite. These exponents are unchanged if the paths
are also confined to horizontal distance $h\log^2h$ from their initial
port.
\end{theorem}

Throughout, $f(h)=h^{a+o(1)}$ means that $f(h)>0$ for all sufficiently
large $h$ and $\log f(h)/\log h\to a$. In particular, the theorem gives
subpower precision; it does not assert bounds by fixed multiples of the
displayed powers. Its last conclusion uses a logarithmically widening
strip, and is distinct from confinement to a fixed multiple of $h$.

The same proof determines a finite free-boundary chord law. In a regular
lattice hexagon of integer side $R$, sum critical weights over paths
joining ordered boundary ports, retaining only paths of diameter at
least $R/10$. Reversal gives a second ordered path and leaves the
normalized law unchanged. Denote the resulting probability by
$\mathbb P_R$.

\begin{theorem}[Hexagon visits]\label{thm:pv:hexagon}
The mass normalizing $\mathbb P_R$ is $R^{3/4+o(1)}$. Uniformly at
vertices in the concentric hexagon of side $R/3$,
\[
 \mathbb P_R(v\in\gamma)=R^{-2/3+o(1)}.
\]
The corresponding upper bound holds at every vertex, and
$\mathbb E_R|\gamma|=R^{4/3+o(1)}$.
\end{theorem}

There is an intermediate strip law in which the distinction between
arches and bridges matters. An \emph{arch} returns to the wall where it
started. From one fixed bottom port, allow both arches and bridges and
condition the path to reach at least $\vartheta h$ layers above the
bottom, where $0<\vartheta<1/2$ is fixed.

\begin{proposition}[Mixed exits]\label{prop:pv:mixed}
The conditioned event has unnormalized mass $h^{-1/4+o(1)}$ and
first-length mass $h^{13/12+o(1)}$. Hence its normalized mean length is
$h^{4/3+o(1)}$. Restricting horizontal distance to $h\log^2h$ preserves
these exponents. The total first-length mass of all exits, without the
height condition, is at most $h^{13/12+o(1)}$.
\end{proposition}

The mixed mean follows directly from a finite vertex-visit identity.
The bridge mean requires an additional argument that transfers mass
from returning arches to crossing paths. It is this argument, rather
than normalization alone, that completes Theorem~\ref{thm:pv:bridge}.
Both theorems concern means in positive critical measures; fixing a
length or prescribing both ports asks a different question.

\subsection{Critical crossing mass and related laws}

Nienhuis's analysis of the dilute $O(n)$ model led to the honeycomb
critical point and the critical-exponent predictions associated with planar
self-avoiding walk~\cite{Nienhuis1982}. Duminil-Copin and Smirnov proved
the connective constant $\sqrt{2+\sqrt2}$ by a parafermionic observable
and its boundary cancellation~\cite{DuminilCopinSmirnov2012}. We use that
local cancellation, prove the particular finite flux identities needed
here, and keep the visited-vertex convention explicit.

The sharper crossing-mass power $h^{-1/4}$ is recorded as a prediction
in \cite[Section~4]{DuminilCopinSmirnov2012}. Beaton, Bousquet-M\'elou,
de Gier, Duminil-Copin and Guttmann proved that the critical bridge
mass tends to zero~\cite{BeatonEtAl2014}. Glazman and
Manolescu~\cite{GlazmanManolescu2017} obtained a quantitative subsequence
bound, as well as invariance of boundary two-point functions in
columnwise rhombic half-planes with Yang--Baxter weights. Krachun and
Panagiotis proved a polynomial upper bound on the critical bridge
mass~\cite{KrachunPanagiotis2026}. These crossing-mass estimates do not
by themselves control a first-length insertion. In particular, the
all-length means considered here are distinct from displacement laws
for uniform walks of a fixed length.

The companion paper \cite[Theorem~1.1]{DiskTransfer2026} proves the same
all-length bridge-mean exponent and retains it inside a fixed corridor
of sufficiently large width $Ch$. It also obtains bridge mass
comparable to $h^{-1/4}$. Our proof of the strip and hexagon results is
given independently below; we do not use the companion's mean theorem.

\subsection{A polygon observable and its construction}

The first-length insertion will be obtained by summing the mass of paths
through a marked vertex. The corresponding algebraic observable counts
surrounding polygons instead of paths. Identify horizontal translates by
$2n$, forming an infinite honeycomb cylinder, and mark two face centers
on one horizontal cut, separated by $n$ sites. Let $F_n(\eta)$ sum over
systems of pairwise vertex-disjoint unoriented simple polygons, each
separating the two marks. A polygon $P$ has weight $\eta x_*^{|P|}$, and
the empty system has weight one. Each polygon crosses the intervening
marked interval, so this is a polynomial in $\eta$; its coefficients are
convergent weighted sums over polygon shapes on the infinite cylinder.

\begin{theorem}[Normalized separator pairing]\label{thm:pv:puncture}
For every compact set $K\subset\{s\in\mathbb C:|\Im s|<\pi\}$ there
is $C_K<\infty$ such that, for all $n\ge1$ and $s\in K$,
\[
 |F_n(2\cosh s)|\le
 C_K n^{1/6+2\Re(s^2)/(3\pi^2)}.
\]
There is a disk about $\eta=0$, independent of $n$, in which $F_n$ has
no zeros. For each fixed $k$, the order-$k$ Taylor coefficient of
$\log F_n$ at zero is $O_k(1+\log n)$. At the positive fugacity two,
\[
 F_n(2)=n^{1/6+o(1)}.
\]
\end{theorem}

To construct $F_n$, record at each of $N=2n$ edges on a cylinder cut
whether the edge is empty or carries either orientation of a path.
The oriented pieces carry local turn phases, and adding a row acts
linearly on the resulting three-state tensors. We construct a common
eigenvector of eigenvalue one whose entries are polynomials in
parameters attached to those edges. Dividing by its empty component
gives the normalized \emph{vacuum}: at the physical honeycomb parameters
its entries are the amplitudes of the half-cylinder path systems.
Pairing two such states with a twist between the marked faces cancels
nonseparating loops and assigns each separating loop weight
$\eta=2\cosh s$.

The normalization $F_n(0)=1$ is essential. Both the finite color formula
for the numerator and the Pfaffian denominator have exponential growth
in $n^2$. The polynomial power survives only after those factors cancel
in the same normalization. Uniformity for complex $s$ also has a
positive-measure consequence: the zero-free disk controls polygon
incompatibilities when the cylinder is opened into the plane.

This construction uses polynomial exchange equations and reductions at
special values of neighboring parameters. Di Francesco and
Zinn-Justin~\cite{DiFrancescoZinnJustin2005} used these ideas for a
multiparameter sum rule in the dense $O(1)$ model. Garbali and
Nienhuis~\cite{GarbaliNienhuis2017Ground,GarbaliNienhuis2017Sum}
developed $q$KZ and fusion recurrences and determinant expressions for
the dilute $O(1)$ model with open boundaries; they also discussed the
polynomial construction at loop weight zero. Here the periodic
twist-$i$ state, its empty-component normalization and its physical
separator pairing are constructed explicitly. The companion's disk
states represent the same homogeneous separator sum;
Section~\ref{c01-c1-the-balanced-puncture-observable} compares the common
scalar polynomial and physical observable. The inhomogeneous vectors
and normalized complex estimates here are proved separately.

The local transfer weights belong to the dilute-loop and three-state
vertex-model setting of Zhou and Batchelor~\cite{ZhouBatchelor1997}.
Related finite recurrences and corner-transfer calculations for
restricted-height dilute models were developed by Warnaar, Pearce,
Seaton and Nienhuis~\cite{WarnaarPearceSeatonNienhuis1994}; Batchelor
and Yung~\cite{BatchelorYung1995} obtained integrable boundary weights
for the ordinary and special transitions of honeycomb polymers.
Bleher and Fokin's six-vertex analysis with domain-wall boundaries
likewise connects determinant partition functions to
orthogonal-polynomial asymptotics~\cite{BleherFokin2006}. These are
methodological comparisons, not inputs for the vacuum or complex-twist
estimate above.

\subsection{Route through the proof}

The proof first obtains two different mass estimates. After the flux
and convergence bounds in Section~\ref{c01-c1-planar-and-cylindrical-setup},
Sections~\ref{c01-c1-the-row-matrices-and-two-fusions}
and~\ref{c01-c1-the-balanced-puncture-observable} construct the polynomial
vacuum and identify its pairing with $F_n$. Differentiating a row at
zero loop fugacity then gives the boundary mass $b_h=h^{-1/4+o(1)}$
in Section~\ref{c01-c1-the-boundary-exponent}. This calculation is
independent of the bulk pairing asymptotics.

Section~\ref{c01-c1-asymptotics-of-the-balanced-observable} proves
Theorem~\ref{thm:pv:puncture}. A finite color sum becomes contour
integrals, with a separate Pfaffian lower bound having bounded
logarithmic error. Periodizing the integrals in an
auxiliary spectral direction permits a circle-determinant expansion.
After that expansion, a positive Fourier decomposition controls the
remaining interacting variables. Their sum stays bounded after its
explicit pressure and twist factors are removed. This spectral period
is distinct from the physical cylinder circumference $2n$.

Section~\ref{c01-c1-localization-of-surrounding-polygons} uses the
zero-free disk and six lattice rotations to localize polygons around
one face: their planar partition function at radius $R$ has exponent
$1/12$. The finite exploration identity in
Section~\ref{c01-c1-the-two-leg-mass} converts this nesting mass into
vertex-visit mass. Summing over $h$ layers and using horizontal
translation gives mixed-exit first-length mass $h^{1+1/12+o(1)}$.
In the hexagon, dividing visit mass by the free-boundary normalizer
$R^{3/4+o(1)}$ gives the visit probability $R^{-2/3+o(1)}$ and hence
the stated mean.

For the strip, the mixed mass could still be carried by returning
arches. Section~\ref{c01-c1-transfer-to-bridges} compares intersecting
bridge--arch pairs with bridge--bridge pairs through a marked tripod.
A translation argument supplies sufficiently many bridge hits on each
deep arch, so bridges retain the same first-length exponent. Only then
do we divide by $b_h$, obtaining
\[
 \mathbb E_h|\gamma|
 =\ell_h/b_h
 =h^{13/12+1/4+o(1)}=h^{4/3+o(1)}.
\]
